\subsection{有效性检验的设计}
\label{sec:diagnostic_design}
实证以Global RRP为固定研究对象，所有比较采用共同评价区间、资产池、252日年化和零无风险利率口径。首先报告五模型的扣费绩效；随后通过消融检验两项惩罚的作用，再结合持仓与事前风险匹配解释风险差异，最后检验成本、执行时点和时间阶段的影响。

消融包括直接持有风险预算参考、将收益短缺惩罚置零及将方差惩罚置零，保留其余估计规则和年度系数。另将评价前选定的首年系数冻结，比较年度更新与初始参数延续。风险匹配将简单基准与货币ETF混合，并用同一历史样本协方差衡量双方当期权重；主模型仍使用原有收缩协方差求解。

成本检验在相同目标路径上改变费率，执行检验将决策顺延一个交易日并重新核算交易。时间阶段固定划分为2018--2020、2021--2023和2024--2026，最后一区间截至样本终点。完整风险匹配方法、估计敏感性、配对区块抽样与收益归因见附录\ref{app:empirical}。
